\section{Experimental program for the hybrid solver}
\label{sec:experiments}

The theory identifies a clear experimental question: does AESP's early
progress compensate for its outer-stage overhead, and can a direct local tail
avoid the expensive part of the active-volume trajectory?  The experiments
must measure this question under a common accuracy and work convention.

\subsection{Hypotheses}

\begin{description}[leftmargin=2.9cm,style=nextline]
\item[H1: early AESP advantage.]
For small \(\alpha\), AESP-LOCAPPR reduces the original PPR error or original
gradient residual more rapidly per degree-weighted operation than standalone
APPR, LocGD, and LocSOR during an initial regime.

\item[H2: diminishing AESP returns.]
The marginal progress per operation decreases after the early regime because
active sets expand, inner tolerances become more demanding, or repeated
outer-stage initialization dominates.

\item[H3: hybrid dominance over AESP.]
A switch to a direct local tail at an adaptive knee reduces total work relative
to the frozen AESP implementation while preserving the same final certificate.

\item[H4: no universal standalone winner.]
AESP, optimal LocSOR, LocCH, LocGD, and APPR have complementary regimes; the
instance-wise best solver depends on graph structure, source, \(\alpha\), and
\(\eps\).
\end{description}

H1 and H2 are motivated by Figure 4 of \citet{huang2025accelerated}: AESP
methods show steep early error reduction, and AESP-LOCAPPR is the most
efficient AESP variant because of online coordinate updates.  H3 is the new
hybrid hypothesis.  H4 motivates using the broader solver set for scientific
calibration without making an instance-wise oracle the official benchmark
target.

\subsection{Methods}

At minimum, compare:
\begin{equation}
\begin{aligned}
    \mathcal M_{\rm exp}:=\{&\AESP\text{-}\LOCAPR,
      \AESP\text{-}\LOCGD,
      \APPR,
      \LOCGD,\\
      \LOCSOR(\omega=1),
      \LOCSOR(\omega_\star),
      \LOCCH,
      \text{hybrid}\}.
\end{aligned}
    \label{eq:experimental-methods}
\end{equation}
The first benchmarked hybrid should be the simplest proof-oriented version in
\cref{alg:hybrid-aesp-locsor}.  Additional variants should isolate one change
at a time:
\begin{enumerate}[leftmargin=2em]
\item fixed outer-stage switch versus objective-bound switch;
\item residual-mass switch versus marginal-work switch;
\item \(\omega=1\) tail versus plain \(\omega_\star\) tail;
\item \(\omega=1\) tail versus safeguarded \(\omega_\star\) tail;
\item FIFO LOCAPPR versus normalized-priority LOCAPR;
\item reset momentum versus any experimental momentum-transfer rule.
\end{enumerate}

\subsection{Primary metrics and correctness}

The primary work measure is
\begin{equation}
    \Work=\sum_{\text{executed }u}d_u,
    \label{eq:experiment-work}
\end{equation}
or the exactly equivalent count of instrumented adjacency-list entries.  Every
method must also report:
\begin{itemize}[leftmargin=2em]
\item outer stages;
\item inner sweeps or coordinate updates;
\item active-set volume by stage and by sweep;
\item queue initialization and stopping-check work;
\item wall-clock time as a secondary engineering metric;
\item peak sparse-state size.
\end{itemize}

A run is valid only if
\begin{equation}
    \norm{D^{-1}(\widehat\pi-\pi)}_\infty\leq\eps,
    \label{eq:experiment-correctness}
\end{equation}
with either an exact small-graph reference or the certified residual condition
\cref{eq:ppr-certificate}.  All baselines must be translated to this same
normalization explicitly.  The inner AESP gradient is never accepted as the
final residual.

\subsection{Instance families}

The suite should mix real and synthetic graphs.

\paragraph{Real graphs.}
Use graphs spanning degree distributions and local community structures,
including the large graphs from the AESP study when practical.  For each
graph, draw sources from degree strata rather than only uniformly: low-degree,
median-degree, and high-degree seeds can produce qualitatively different
active trajectories.

\paragraph{Synthetic mechanisms.}
Include the following controlled families:
\begin{itemize}[leftmargin=2em]
\item center- and leaf-seeded stars, which expose high-degree boundary costs;
\item paths and long spiders, which expose diffusive propagation and local
      spectral obstructions;
\item core--boundary--exterior graphs of varying boundary volume, following
      the mechanism in \citet{fountoulakis2026complexity};
\item expanders and dense local cores, which test whether the hybrid switches
      too early when global mixing is favorable;
\item graphs with the same local core but different remote graph size, which
      test graph-size independence.
\end{itemize}

\subsection{Parameter grid and reporting}

Use a logarithmic grid in \(\alpha\), emphasizing
\(
10^{-3}\) to \(10^{-1}
\), and several \(\eps\) values spanning coarse and high-accuracy regimes.
For every run, record
\begin{equation}
   (G,s,\alpha,\eps,\text{seed},\text{ordering},\text{switch rule},
    \text{stopping rule},\text{code version}).
    \label{eq:experiment-metadata}
\end{equation}
Plots should show error versus cumulative work and active volume versus work,
not only final bar charts.  The hybrid claim is specifically about the shape
of the trajectory and the location of the switching knee.

\subsection{Falsification criteria}

The hybrid hypothesis should be considered falsified, or at least materially
weakened, if any of the following persists after implementation debugging:
\begin{enumerate}[leftmargin=2em]
\item the reference hybrid fails to beat frozen AESP on a diverse held-out
      suite by more than measurement noise;
\item all gains are explained by a stopping-rule mismatch;
\item a standalone existing solver already passes the proposed benchmark
      threshold on almost every hidden instance;
\item the switch requires graph-family labels or hand-tuned per-instance
      parameters unavailable to an agent;
\item queue initialization or residual reconstruction makes the nominal local
      work advantage disappear under exact accounting.
\end{enumerate}
